% year: תשע"ח
% type: מבחן
% semester: ב
% moed: א
% number: 3ב
% points: 10
% categories: taylor, derivatives
% source: exams/מבחנים/תשעח/מועד א תשעח סמסטר ב.pdf

\begin{question}
רשמו פולינום מקלורן מסדר 2 עבור הפונקציה $y(x)$ אשר מוגדרת בצורה סתומה על ידי המשוואה $y\cdot\sin x-x\cdot\cos x+y=2$.
\end{question}

\begin{hint}
הציבו $x=0$ כדי למצוא את $y(0)$, ואז גזרו את המשוואה (גזירה סתומה) פעמיים והציבו $x=0$.
\end{hint}

\begin{solution}
\textbf{ערך ב-$0$:} בהצבה $x=0$: $y(0)\cdot0-0+y(0)=2$, ולכן $y(0)=2$.

\textbf{נגזרת ראשונה} (גזירה סתומה של שני האגפים לפי $x$, כאשר $y=y(x)$):
\[
y'\sin x+y\cos x-\cos x+x\sin x+y'=0.\tag{*}
\]
בהצבה $x=0$: $y(0)-1+y'(0)=0$, ולכן $y'(0)=1-2=-1$.

\textbf{נגזרת שנייה} (גזירה של (*)):
\[
y''\sin x+y'\cos x+y'\cos x-y\sin x+\sin x+\sin x+x\cos x+y''=0.
\]
בהצבה $x=0$: $2y'(0)+y''(0)=0$, ולכן $y''(0)=2$.

\textbf{פולינום מקלורן מסדר 2:}
\[
P_2(x)=y(0)+y'(0)x+\frac{y''(0)}{2}x^2=2-x+x^2.
\]
(בדיקה: מהמשוואה $y=\frac{2+x\cos x}{1+\sin x}$, ופיתוח ישיר נותן אותו פולינום.)
\end{solution}
